\documentclass{article}

\input{~/University/hwpreamble.tex}

\begin{document}

\title{Topología - Ejercicios profesor Ruíz}
\author{Victoria Eugenia Torroja Rubio}
\date{7/9/26}

\maketitle

\begin{ej}[Número 0.8]
Sea $\displaystyle X $ un espacio compacto y sea $\displaystyle Y = \mathcal{C}\left(X, \R^{n}\right) $ el conjunto de las funcioens continuas $\displaystyle X \to \R^{n} $. Demostrar que $\displaystyle d\left(f,g\right)=\max \|f-g\| $ es una distancia completa en $\displaystyle Y $.
\end{ej}

\begin{sol}
Entendemos que $\displaystyle X $ es un espacio compacto si y solo si todo recubrimiento abierto admite un subrecubrimiento finito. \\ \\
Veamos primero que $\displaystyle d\left(f,g\right)=\max \|f-g\| $ es una distancia en $\displaystyle Y $.
\begin{itemize}
	\item En primer lugar, tenemos que ver que $\displaystyle d : X \times X \to \R $ está bien definida, es decir, que cada función $\displaystyle f \in Y $ alcanza su máximo en $\displaystyle X $. Veamos que si $\displaystyle f\in Y $ y $\displaystyle X $ es compacto, $\displaystyle f\left(X\right) $ es compacto. Sea $\displaystyle \left\{ U_{i}\right\} _{i \in I} $ es un recubrimiento abierto de $\displaystyle f\left(X\right) $. Como $\displaystyle f $ es continua se tiene que $\displaystyle f^{-1}\left(U_{i}\right) $ es abierto $\displaystyle \forall i \in I $. Así, 
		\[f\left(X\right)\subset \bigcup_{i \in I}U_{i} \Rightarrow X \subset f^{-1}\left(\bigcup_{i \in I}U_{i}\right)=\bigcup_{i \in I}f^{-1}\left(U_{i}\right) .\]
		Por tanto, $\displaystyle \left\{ f^{-1}\left(U_{i}\right)\right\} _{i \in I} $ es un recubrimiento abierto de $\displaystyle X $. Como $\displaystyle X $ es compacto, existe $\displaystyle J \subset I $ finito tal que $\displaystyle \left\{ f^{-1}\left(U_{j}\right)\right\} _{j\in J} $ es un recubrimiento de $\displaystyle X $, por lo que debe ser que $\displaystyle \left\{ U_{j}\right\} _{j \in J} $ es un recubrimiento finito de $\displaystyle f\left(X\right) $. En efecto,
		\[f\left(X\right) \subset f\left(\bigcup_{j \in J}f^{-1}\left(U_{j}\right)\right)=\bigcup_{j \in J}f\left(f^{-1}\left(U_{j}\right)\right)=\bigcup_{j \in J}U_{j} .\]
		Como $\displaystyle f\left(X\right)\subset \R^{n} $ es compacto, tendremos que es cerrado y acotado. En concreto, la función $\displaystyle h = \|f-g\| : X \to \R $ es continua por lo que $\displaystyle h\left(X\right) $ es un compacto de $\displaystyle \R $. Así, existen $\displaystyle \alpha = \inf h\left(X\right) $ y $\displaystyle \beta = \sup h\left(X\right) $ y como $\displaystyle h\left(X\right) $ es cerrado y acotado se tiene que 
		\[\alpha, \beta \in \overline{h\left(X\right)}=h\left(X\right) .\]
		Por tanto, la función $\displaystyle h = \|f-g\| $ alcanza su máximo en $\displaystyle X $ y tenemos que la función distancia dada en el enunciado está bien definida.
	\item Veamos que se trata de una función distancia. 
		\begin{enumerate}
		\item En primer lugar, es sencillo ver que como $\displaystyle \|f-g\| \geq 0 $, tendremos que $\displaystyle d\left(f,g\right) \geq 0 $.
		\item Si $\displaystyle f \equiv g $ claramente $\displaystyle d\left(f,g\right)=0 $. Recíprocamente, si $\displaystyle d\left(f,g\right)=0 $ y existe $\displaystyle x_{0} \in X $ con $\displaystyle f\left(x_{0}\right)\neq g\left(x_{0}\right) $, tendremos que $\displaystyle \max \|f\left(x\right)-g\left(x\right)\| > 0 $ por lo que necesariamente debe ser que $\displaystyle f \equiv g $.
		\item Es trivial que la propiedad simétrica se cumple.
		\item Veamos que se cumple la propiedad triangular. Si $\displaystyle f,g,h \in Y $,
			\[
			\begin{split}
				d\left(f,g\right)= &\max \|f-g\|=\max \|f-h+h-g\|\leq \max \left(\|f-h\| + \|h-g\|\right) \\
				\leq& \max \|f-h\| + \max \|h - g\| = d\left(f,h\right) + d\left(h,g\right) .
			\end{split}
			\]
		\end{enumerate}
	\item Veamos que es una distancia completa. Sea $\displaystyle \left\{ f_{m}\right\} _{m\in\N}\subset Y $ una sucesión de Cauchy. Tenemos que para $\displaystyle \epsilon > 0 $ existe $\displaystyle n_{0} \in \N $ tal que si $\displaystyle m,k\geq n_{0} $,
		\[d\left(f_{m}-f_{k}\right)=\max \|f_{m}-f_{k}\|<\epsilon \iff \|f_{m}\left(x\right)-f_{k}\left(x\right)\|<\epsilon, \; \forall x \in X .\]
		Podemos ver que para cada $\displaystyle x_{0} \in X $ tenemos que $\displaystyle \left\{ f_{m}\left(x_{0}\right)\right\} _{m \in \N} \subset \R^{n} $ también es sucesión de Cauchy. Como $\displaystyle \R^{n} $ es completo, existe 
		\[f\left(x_{0}\right):=\lim_{m \to \infty}f_{m}\left(x_{0}\right) \in \R^{n} .\]
		Así, definimos la función $\displaystyle f : X \to \R^{n} $ definida como $\displaystyle f\left(x\right)=\lim_{m \to \infty}f_{m}\left(x\right) $. Veamos que la sucesión $\displaystyle \left\{ f_{m}\right\} _{m\in\N} $ converge a $\displaystyle f $. Dado $\displaystyle \epsilon > 0 $, existe $\displaystyle n_{0} \in \N $ tal que si $\displaystyle m,k \geq n_{0} $ entonces
		\[\|f_{m}\left(x\right)-f_{k}\left(x\right)\|<\epsilon, \; \forall x \in X .\]
		Tomando el límite cuando $\displaystyle m \to \infty $ en la desigualdad anterior obtenemos que
		\[\|f\left(x\right)-f_{k}\left(x\right)\|<\epsilon  .\]
		Por tanto, $\displaystyle f_{m} \to f $. Finalmente, veamos que $\displaystyle f \in Y $. Sea $\displaystyle V \subset \R^{n} $ un abierto, queremos ver que $\displaystyle f^{-1}\left(V\right) $ es abierto. Sea $\displaystyle x_{0} \in f^{-1}\left(V\right) $, entonces $\displaystyle f\left(x_{0}\right)\in V $ y como $\displaystyle V $ es abierto existe $\displaystyle \epsilon > 0 $ tal que $\displaystyle B\left(f\left(x_{0}\right), \epsilon \right)\subset V $. 
	Como $\displaystyle f_{m} \to f $, podemos escoger $\displaystyle m \in \N $ suficientemente grande tal que $\displaystyle \|f_{m}\left(x\right)-f\left(x\right)\|<\frac{\epsilon }{3} $, $\displaystyle \forall x \in X $, en particular
	\[\|f_{m}\left(x_{0}\right)-f\left(x_{0}\right)\|<\frac{\epsilon }{3} .\]
Como $\displaystyle f_{m} $ es continua, $\displaystyle f^{-1}\left(B\left(f_{m}\left(x_{0}\right), \frac{\epsilon }{3}\right)\right) $ es un abierto en $\displaystyle X $. Definimos
\[U :=f^{-1}\left(B\left(f_{m}\left(x_{0}\right), \frac{\epsilon }{3}\right)\right) .\]
Tenemos que $\displaystyle x_{0} \in  U $. Si $\displaystyle x \in U $ tenemos que $\displaystyle \|f_{m}\left(x\right)-f_{m}\left(x_{0}\right)\|<\frac{\epsilon }{3} $, por lo que 
\[\|f\left(x\right)-f\left(x_{0}\right)\| \leq \|f\left(x\right)-f_{m}\left(x\right)\| + \|f_{m}\left(x\right)-f_{m}\left(x_{0}\right)\| + \|f_{m}\left(x_{0}\right)-f\left(x_{0}\right)\| < \frac{\epsilon }{3} + \frac{\epsilon }{3} + \frac{\epsilon }{3}=\epsilon  .\]
Por tanto, $\displaystyle f\left(x\right)\in B\left(f\left(x_{0}\right), \epsilon \right)\subset V $. Así, $\displaystyle U \subset f^{-1}\left(V\right) $ y tenemos que $\displaystyle f^{-1}\left(V\right) $ es abierto. 
\end{itemize}
\end{sol}

\begin{ej}[Número 1.6]
	\begin{enumerate}
	\item Sean $\displaystyle \mathcal{T}_{1} $ y $\displaystyle \mathcal{T}_{2} $ dos topologías en un conjunto $\displaystyle X $. Mostrar que su intersección $\displaystyle \mathcal{T}_{1} \cap \mathcal{T}_{2} $ es también una topología. ¿Se puede decir lo mismo de su unión $\displaystyle \mathcal{T}_{1} \cup \mathcal{T}_{2} $? Construir la topología menos fina que contiene a esa unión.
	\item Estudiar el caso de intersecciones y uniones arbitrarias de topologías.
	\item La menor topología que contiene a una colección cualquiera de subconjuntos de $\displaystyle X $ se denomina topología generada por la colección. Describirla. ¿Qué tiene esto que ver con el concepto de base de una topología?
	\end{enumerate}
\end{ej}
\begin{sol}
Resolvemos los distintos apartados.
\begin{enumerate}
\item Veamos que $\displaystyle \mathcal{T}_{1}\cap \mathcal{T}_{2} $ es una topología en $\displaystyle X $. 
	\begin{itemize}
	\item Como $\displaystyle \mathcal{T}_{1} $ y $\displaystyle \mathcal{T}_{2} $ son topologías en $\displaystyle X $, por definición se sigue que $\displaystyle \emptyset, X \in \mathcal{T}_{1} \cap \mathcal{T}_{2} $. 
	\item Si $\displaystyle \left\{ U_{i}\right\} _{i \in I}\subset \mathcal{T}_{1}\cap \mathcal{T}_{2} $, tenemos que por ser $\displaystyle \mathcal{T}_{1} $ y $\displaystyle \mathcal{T}_{2} $ topologías, 
		\[\bigcup_{i \in I}U_{i} \in \mathcal{T}_{1} \quad \text{y} \quad \bigcup_{i \in I}U_{i} \in \mathcal{T}_{2} ,\]
		por lo que $\displaystyle \bigcup_{i \in I}U_{i} \in \mathcal{T}_{1} \cap \mathcal{T}_{2} $. 
	\item Supongamos que $\displaystyle U_{1}, U_{2}\in \mathcal{T}_{1} \cap \mathcal{T}_{2} $. Tenemos que por ser $\displaystyle \mathcal{T}_{1} $ y $\displaystyle \mathcal{T}_{2} $ topologías, $\displaystyle U_{1} \cap U_{2}\in \mathcal{T}_{j} $ con $\displaystyle j \in \left\{ 1, 2\right\}  $, por lo que $\displaystyle U_{1} \cap U_{2} \in \mathcal{T}_{1} \cap \mathcal{T}_{2} $. 
	\end{itemize}
En general, no es cierto que $\displaystyle \mathcal{T}_{1} \cup \mathcal{T}_{2} $ sea una topología. Consideremos un conjunto $\displaystyle X $ infinito y las topologías $\displaystyle \mathcal{T}_{1}=\mathcal{T}_{a} $ para cierto $\displaystyle a \in X $, y $\displaystyle \mathcal{T}_{2}=\mathcal{T}_{CF} $. Dado $\displaystyle x \in X $ con $\displaystyle x \neq a $ podemos considerar los conjuntos 
\[ A = \left\{ a,x\right\} \in \mathcal{T}_{1} \quad \text{y} \quad B = X \backslash \left\{ a, x_{1}, \ldots, x_{r}\right\} \in \mathcal{T}_{2} .\]
Tendremos que $\displaystyle A \cap B = \left\{ x\right\}   $, que no pertenece a $\displaystyle \mathcal{T}_{1} $ ni a $\displaystyle \mathcal{T}_{2} $. Por tanto, $\displaystyle \mathcal{T}_{1} \cup \mathcal{T}_{2} $ no es una topología en general. La topología más fina que contiene a $\displaystyle \mathcal{T}_{1}\cup \mathcal{T}_{2} $ es la que la tiene por subbase. Denotemos por $\displaystyle \mathcal{T}_{\mathcal{T}_{1}\cup\mathcal{T}_{2}} $ 
a la topología generada por la subbase $\displaystyle \mathcal{T}_{1} \cup \mathcal{T}_{2} $ y sea $\displaystyle \mathcal{T} $ una topología que contiene a $\displaystyle \mathcal{T}_{1} \cup \mathcal{T}_{2} $. 
Si $\displaystyle U \in \mathcal{T}_{\mathcal{T}_{1}\cup\mathcal{T}_{2}} $, tenemos que $\displaystyle U $ es unión de intersecciones finitas de elementos de $\displaystyle \mathcal{T}_{1} \cup \mathcal{T}_{2} $, esto es
\[U = \bigcup_{i \in I}\left(\bigcap_{j = 1}^{r_{i}}V_{j}\right) , \quad V_{j}\in \mathcal{T}_{1} \cup \mathcal{T}_{2},\]
Claramente tendremos que $\displaystyle U \in \mathcal{T} $ por ser esta una topología que contiene a $\displaystyle \mathcal{T}_{1} \cup \mathcal{T}_{2} $. Por tanto, hemos visto que $\displaystyle \mathcal{T}_{\mathcal{T}_{1} \cup \mathcal{T}_{2}} $ es la topología más fina que contiene a $\displaystyle \mathcal{T}_{1} \cup \mathcal{T}_{2} $. 
\item La intersección arbitraria de topologías es una topología, en efecto, consideremos la familia $\displaystyle \left\{ \mathcal{T}_{i}\right\} _{i \in I} $ de topologías en $\displaystyle X $:
	\begin{itemize}
	\item Como $\displaystyle \emptyset, X \in \mathcal{T}_{i} $, $\displaystyle \forall i \in I $, tenemos que $\displaystyle \emptyset, X \in \bigcap_{ i \in I}\mathcal{T}_{i} $. 
	\item Dado $\displaystyle \left\{ U_{j}\right\} _{j \in J} \subset \bigcap_{i \in I}\mathcal{T}_{i}$, tenemos que $\displaystyle \bigcup_{j \in J}U_{j}\in \mathcal{T}_{i} $, $\displaystyle \forall i \in I $, por lo que $\displaystyle \bigcup_{j \in J}U_{j}\in \bigcap_{i \in I}\mathcal{T}_{i} $. 
	\item Si $\displaystyle U_{1} , U_{2} \in \bigcap_{i \in I}\mathcal{T}_{i} $, entonces $\displaystyle U_{1} \cap U_{2} \in \mathcal{T}_{i} $, $\displaystyle \forall i \in I $, por lo que $\displaystyle U_{1} \cap U_{2} \in \bigcap_{i \in I}\mathcal{T}_{i} $. 
	\end{itemize}
	Antes vimos que la unión de topologías en general no es topología, por lo que lo mismo va a suceder con la unión arbitraria de topologías. 
\item La topología generada por una colección es la topología generada por la subbase que es esa colección. De esta forma, el conjunto de intersecciones finitas de la colección dada nos da una base de una topología, que es precisamente la topología más fina que contiene a la colección. \footnote{Esto es siempre y cuando $\displaystyle X $ se pueda expresar como unión de elementos de la colección.} 
\end{enumerate}
\end{sol}

\end{document}
